% ============================================================
%  Mission 50' — Nombres complexes
%  Canevas de restitution (1re STI2D) — Équipe Écran + bilan
%  Compilation : pdflatex (TeXstudio : F5 ou Compilation rapide)
% ============================================================
\documentclass[11pt,a4paper]{article}

\usepackage[utf8]{inputenc}
\usepackage[T1]{fontenc}
\usepackage[french]{babel}
\usepackage{lmodern}
\usepackage{amsmath,amssymb}
\usepackage[margin=1.8cm]{geometry}
\usepackage{tabularx}
\usepackage{enumitem}
\usepackage{xcolor}
\usepackage{array}
\usepackage[table]{xcolor}   % au lieu de : \usepackage{xcolor}

% ---------- Réglages ----------
\setlength{\parindent}{0pt}
\setlength{\parskip}{5pt}
\renewcommand{\arraystretch}{1.4}

% Case à cocher
\newcommand{\case}{$\square$}

% Ligne de réponse à compléter
\newcommand{\pointilles}{\makebox[3.5cm]{\dotfill}}
\newcommand{\pointilleslarge}{\makebox[0.92\textwidth]{\dotfill}}

% En-tête de feuille : #1 = numéro (ex. 1/4), #2 = titre, #3 = badge équipe
\newcommand{\feuille}[3]{%
\begin{center}
  {\bfseries\large Feuille #1 --- #2}\\[2pt]
  {\small\colorbox{gray!15}{\parbox{0.5\textwidth}{\centering #3}}}
\end{center}}

\begin{document}

% ============================================================
%  FEUILLE 1/4 — COMPTE RENDU DE MISSION
% ============================================================
\feuille{1/4}{Compte rendu de mission}{ÉQUIPE ÉCRAN · 1 PAR BINÔME}

Mission 50' --- les nombres complexes · 1\up{re} STI2D. Ce compte rendu est le
rendu écrit de votre binôme : soignez la rédaction, la mise en commun se fera
en classe.

\begin{center}
\renewcommand{\arraystretch}{2}
\begin{tabularx}{\textwidth}{|l|X|l|X|l|X|l|X|}
  \hline
  \textbf{Binôme n°} & & \textbf{Pilote} & & \textbf{Rapporteur} & &
  \textbf{Graine n°} & \\ \hline
\end{tabularx}
\end{center}

Rôles échangés à la phase 2 --- écrivez les deux noms, le n° de graine affiché
en haut de votre fichier Écran.

\subsection*{A --- Phase 1 : Nos deux missions de recherche (chapitres 1 \& 2)}

\textbf{Mission 1 choisie} : \pointilles\ (recopiez son titre)

\textbf{Notre réponse en 1 ou 2 phrases :}

\pointilles\par\pointilles

\textbf{Notre source (site, vidéo, article) :}

\pointilles

\vspace{6pt}

\textbf{Mission 2 choisie} : \pointilles\ (recopiez son titre)

\textbf{Notre réponse en 1 ou 2 phrases :}

\pointilles\par\pointilles

\textbf{Notre source (site, vidéo, article) :}

\pointilles

\subsection*{B --- Phase 2 : Nos exercices (chapitres 3 \& 4)}

\begin{center}
\renewcommand{\arraystretch}{2.2}
\begin{tabularx}{\textwidth}{|X|X|c|}
  \hline
  \textbf{Exercice} & \textbf{Notre réponse rédigée (forme attendue :}
  & \textbf{Sûrs ?} \\
   & $z = a + ib$, $|z| = \dots$, $z = r(\cos\theta + i\sin\theta)$\dots)
  & $\checkmark$ / $\pm$ / $\times$ \\ \hline
  Ex 1 --- forme algébrique & & \\ \hline
  Ex 2 --- forme algébrique & & \\ \hline
  Ex 3 --- forme algébrique & & \\ \hline
  Ex 4 --- forme algébrique & & \\ \hline
  Ex 1 --- forme trigonométrique & & \\ \hline
  Ex 2 --- forme trigonométrique & & \\ \hline
  Ex 3 --- forme trigonométrique & & \\ \hline
  Ex 4 --- forme trigonométrique & & \\ \hline
\end{tabularx}
\end{center}

\subsection*{C --- Labo : le cycle caché des puissances de $i$}

Recopiez la ligne du bas du tableau du chapitre 3 ($n = 1$ à 12) :

\begin{center}
\renewcommand{\arraystretch}{2}
\begin{tabular}{|c|*{12}{c|}}
  \hline
  $n$ & 1 & 2 & 3 & 4 & 5 & 6 & 7 & 8 & 9 & 10 & 11 & 12 \\ \hline
  $i^{n}$ & & & & & & & & & & & & \\ \hline
\end{tabular}
\end{center}

\textbf{Défi :} $i$ puissance \pointilles\ = \pointilles\
(réponse : $1$, $i$, $-1$ ou $-i$) --- justification par la division
par 4 :

\pointilles\par\pointilles

\vspace{4pt}
\emph{Aucune correction n'est fournie dans les fichiers élèves : la mise en
commun des réponses se fait en classe, phase 4.}

\newpage
% ============================================================
%  FEUILLE 2/4 — QUIZ
% ============================================================
\feuille{2/4}{Quiz --- nos 10 réponses}{ÉQUIPE ÉCRAN · 1 PAR BINÔME}

Une seule tentative par question, aucun score ni correction affichés :
reportez chaque réponse ici dès que le binôme a cliqué. La mise en commun se
fera en classe.

\begin{center}
\renewcommand{\arraystretch}{2.4}
\begin{tabular}{|c|p{5.5cm}|c|p{5.5cm}|}
  \hline
  \textbf{N°} & \textbf{Notre réponse (la lettre ET la valeur,
  ex. « B : $-1$ »)} & \textbf{N°} & \textbf{Notre réponse} \\ \hline
  1 & & 6 & \\ \hline
  2 & & 7 & \\ \hline
  3 & & 8 & \\ \hline
  4 & & 9 & \\ \hline
  5 & & 10 & \\ \hline
\end{tabular}
\end{center}

\textbf{La question qui nous a le plus résisté :} n° \pointilles\ ---
pourquoi ?

\pointilles\par\pointilles

\textbf{Une question d'observation à garder pour la classe} (le cycle de $i$,
un laboratoire, une application\dots) :

\pointilles\par\pointilles

\vfill
\begin{center}
\footnotesize Équipe Écran · Mission 50' --- les nombres complexes ·
1\up{re} STI2D · graine n° : \pointilles
\end{center}

\newpage
% ============================================================
%  FEUILLE 3/4 — CARTES DE TRANSMISSION
% ============================================================
\feuille{3/4}{Nos 3 cartes de transmission}{ÉQUIPE ÉCRAN · 1 PAR BINÔME}

Trois cartes à présenter oralement ($\approx$ 1 minute par carte) à l'équipe
Papier lors de la passerelle, phase 4. Elles seront croisées avec les cartes
« papier » (fiche F). Une carte = trois cases : \textbf{la découverte, la
preuve, à quoi ça sert}.

\newcommand{\carte}[1]{%
\begin{center}
\renewcommand{\arraystretch}{1.9}
\begin{tabularx}{\textwidth}{|c|X|}
  \hline
  \multicolumn{2}{|l|}{\cellcolor{gray!15}\textbf{CARTE #1}} \\ \hline
  \textcircled{1} & \textbf{La découverte} --- une phrase courte
  \\[1.1cm] \hline
  \textcircled{2} & \textbf{La preuve} --- un calcul, une valeur ou une figure
  \\[1.1cm] \hline
  \textcircled{3} & \textbf{Ça sert à quoi ?} --- un domaine technique concret
  \\[1.1cm] \hline
\end{tabularx}
\end{center}}

\carte{1}
\carte{2}
\carte{3}

\vfill
\begin{center}
Binôme : \pointilles\ · Pilote : \pointilles\ · Rapporteur : \pointilles
\end{center}

\newpage
% ============================================================
%  FEUILLE 4/4 — FICHE BILAN INDIVIDUELLE
% ============================================================
\feuille{4/4}{Ma fiche bilan individuelle}{LES DEUX ÉQUIPES · 1 PAR ÉLÈVE}

\begin{center}
\renewcommand{\arraystretch}{2}
\begin{tabularx}{\textwidth}{|l|X|l|c|l|X|}
  \hline
  \textbf{Nom, prénom} & & \textbf{Équipe} & \case{} Écran \; \case{} Papier &
  \textbf{Binôme n°} & \\ \hline
\end{tabularx}
\end{center}

À conserver précieusement : le prochain cours s'appuiera sur cette fiche. La
correction des quiz et exercices se fera en classe.

\section*{1 --- Je retiens (à compléter)}

\begin{itemize}
  \item $i$ est le nombre tel que $i^2 = \pointilles$ ;
  \item $z = a + ib$ : $a$ s'appelle la partie \pointilles, $b$ la partie
        \pointilles ;
  \item Dans le plan, le complexe $z$ est le point $M(a\,;\,b)$. Le module
        $|z| = \sqrt{\pointilles}$ est une \pointilles\ (distance $OM$),
        l'argument est une \pointilles\ (angle) ;
  \item $z = r(\cos\theta + i\sin\theta)$ : $r$ est le \pointilles,
        $\theta$ l'\pointilles
\end{itemize}

\section*{2 --- Les 3 questions que j'apporte au cours}

1. \pointilleslarge

2. \pointilleslarge

3. \pointilleslarge

\section*{3 --- Où j'en suis (cocher une case par ligne)}

\begin{center}
\renewcommand{\arraystretch}{2}
\begin{tabularx}{\textwidth}{|X|c|c|c|}
  \hline
  \textbf{Objectif de la mission} & \textbf{J'y suis $\checkmark$}
  & \textbf{Presque} & \textbf{Pas encore} \\ \hline
  Connaître l'histoire : pourquoi on a « inventé » ces nombres, et qui les a
  apprivoisés & \case{} & \case{} & \case{} \\ \hline
  Citer au moins trois domaines techniques où les complexes sont
  indispensables & \case{} & \case{} & \case{} \\ \hline
  Écrire et calculer avec $z = a + ib$ (forme algébrique)
  & \case{} & \case{} & \case{} \\ \hline
  Placer un nombre complexe dans le plan, y lire module et argument
  & \case{} & \case{} & \case{} \\ \hline
  Passer de la forme algébrique à la forme $z = r(\cos\theta + i\sin\theta)$,
  et inversement & \case{} & \case{} & \case{} \\ \hline
\end{tabularx}
\end{center}

\section*{4 --- Ma preuve préférée}

Le calcul, la figure ou le fait de la mission dont je suis le plus fier·e
(recopiez-le, il servira au prochain cours) :

\pointilles\par\pointilles\par\pointilles

\end{document}
